\section{Agent IA de traitement des chèques}

Cette section présente l'approche utilisée pour le traitement automatisé
des chèques bancaires. Contrairement aux documents KYC (CIN, justificatifs
de domicile) dont l'information recherchée est de même nature, le chèque
mêle sur un même support \textbf{trois types d'information visuellement très
différents} :
\begin{itemize}
    \item des \textbf{champs manuscrits} (montant en chiffres, montant en
    lettres, nom du bénéficiaire), rédigés à la main par l'émetteur ;
    \item des \textbf{zones imprimées} normalisées (numéro de compte, bande
    magnétique CMC7) ;
    \item une \textbf{zone graphique} (la signature), dont seule la présence
    doit être vérifiée.
\end{itemize}

Cette hétérogénéité interdit l'usage d'un modèle unique. Le système repose
donc sur une \textbf{architecture hybride combinant trois modèles
complémentaires}, chacun spécialisé sur le type d'information qu'il sait
traiter au mieux.

\subsection*{Objectif et architecture globale}

\subsubsection{Objectif}

L'objectif de cette partie est de développer un système intelligent capable
d'automatiser la lecture et la validation d'un chèque bancaire marocain afin
d'alimenter le processus de back-office de manière automatisée.

Les champs traités sont les suivants :
\begin{itemize}
    \item Le montant en chiffres.
    \item Le montant en lettres.
    \item Le nom du bénéficiaire.
    \item Le numéro de compte.
    \item La bande CMC7 (ligne magnétique de bas de chèque).
    \item La présence de la signature.
\end{itemize}

\subsubsection{Architecture hybride à trois modèles}

Le traitement d'un chèque repose sur trois modèles d'intelligence
artificielle complémentaires, orchestrés au sein du module
\texttt{tools/extractor.py} :

\begin{itemize}
    \item \textbf{YOLOv8} pour la \textit{détection des zones} structurées
    du chèque (numéro de compte, signature, bande CMC7) ;
    \item \textbf{TrOCR} (variante \textit{printed}) pour la reconnaissance
    optique du \textit{texte imprimé} des zones détectées (numéro de compte
    et CMC7) ;
    \item \textbf{Qwen2-VL-2B} affiné par \textit{LoRA} pour la lecture des
    \textit{champs manuscrits} (montants et bénéficiaire), où l'OCR
    classique échoue.
\end{itemize}

Le processus global de traitement d'un chèque se décompose comme suit :
\begin{enumerate}
    \item \textbf{Lecture du manuscrit} : l'image complète du chèque est
    transmise à Qwen2-VL, qui renvoie directement un JSON structuré
    (montant en chiffres, montant en lettres, bénéficiaire).
    \item \textbf{Détection des zones} : YOLOv8 localise les trois zones
    (numéro de compte, signature, CMC7) et en découpe les régions
    (\textit{crops}).
    \item \textbf{OCR des zones imprimées} : les régions correspondant au
    numéro de compte et à la bande CMC7 sont transmises à TrOCR pour
    l'extraction du texte.
    \item \textbf{Validation métier} : les données extraites par les deux
    pipelines sont confrontées à un ensemble de règles produisant une
    décision structurée.
\end{enumerate}

Ce découplage permet d'attribuer chaque tâche au modèle le plus adapté :
un modèle Vision-Langage pour le raisonnement sémantique sur l'écriture
manuscrite, et un pipeline de détection + OCR léger et rapide pour les
zones imprimées au gabarit régulier.

% ─────────────────────────────────────────────────────────────────────────────
\subsection*{Détection des zones par YOLOv8}

\subsubsection{Objectif et classes détectées}

La première brique du système est un modèle YOLOv8 chargé de localiser les
zones structurées du chèque. Trois classes ont été définies :

\begin{itemize}
    \item \texttt{num\_compte} — le bloc imprimé contenant le numéro de
    compte du tireur ;
    \item \texttt{signature} — la zone de signature, en bas à droite du
    chèque ;
    \item \texttt{cmc7} — la ligne magnétique normalisée (CMC7) figurant en
    bas du chèque, qui encode le numéro de chèque, le code agence et le
    numéro de compte.
\end{itemize}

Le choix de YOLOv8 pour ces zones repose sur la même logique que pour la
CIN : ces champs occupent une position \textbf{stable et géométriquement
régulière} sur le gabarit du chèque, configuration pour laquelle un
détecteur basé sur la localisation spatiale est particulièrement performant.

\subsubsection{Génération du dataset synthétique}

À l'image du pipeline CIN, un dataset synthétique a été généré
automatiquement à partir de templates de chèques représentatifs des
principales banques marocaines (Attijariwafa Bank, BMCE, BCP, CIH, etc.).
Pour chaque image générée :
\begin{itemize}
    \item des identités, montants et numéros aléatoires ont été générés ;
    \item les champs ont été positionnés dynamiquement sur les templates ;
    \item un décalage spatial aléatoire a été appliqué pour simuler la
    variabilité de numérisation (scan décadré, photo prise à la main) ;
    \item les \textit{bounding boxes} YOLO des trois classes ont été
    calculées automatiquement.
\end{itemize}

Le dataset a été réparti entre les ensembles d'entraînement, de validation
et de test selon le même découpage que pour la CIN.

\subsubsection{Entraînement du modèle YOLOv8}

Comme pour la CIN, le modèle \textbf{YOLOv8n} (variante nano, $\approx$ 3,2
millions de paramètres) a été retenu pour son compromis favorable entre
légèreté et précision sur des objets structurés. L'entraînement a été
réalisé par \textbf{transfert d'apprentissage} à partir des poids
pré-entraînés sur le dataset MS COCO, accélérant la convergence et
réduisant le risque de surapprentissage sur un dataset de taille modeste.

\begin{table}[H]
\centering
\label{tab:yolo_cheque_training}
\begin{tabular}{|l|c|}
\hline
\textbf{Paramètre} & \textbf{Valeur} \\
\hline
Modèle de base        & YOLOv8n (pré-entraîné COCO) \\
Nombre de classes     & 3 (num\_compte, signature, cmc7) \\
Nombre d'époques      & 50 \\
Taille des images     & 640 $\times$ 640 \\
Batch size            & 8 \\
Mode AMP              & Activé \\
Patience (early stopping) & 15 \\
\hline
\end{tabular}
\caption{Configuration d'entraînement du modèle YOLOv8n pour les chèques}
\end{table}

Le meilleur modèle observé en validation (maximisant le mAP@0.5) est
sauvegardé sous le nom \texttt{best.pt}, puis utilisé en production
(\texttt{modele\_final/best.pt}).

\subsubsection{Optimisation de l'inférence}

Plusieurs optimisations ont été mises en place lors de la phase de
détection (\texttt{extract\_zones}) afin de garantir des temps de réponse
faibles et des détections fiables :

\begin{itemize}
    \item \textbf{Suppression des non-maxima (NMS)} : un seuil
    $\text{IoU} = 0{,}4$ est appliqué pour éviter la double détection d'une
    même zone. Pour chaque classe, seule la détection de plus forte
    confiance est conservée.
    \item \textbf{Découpe en mémoire} : les \textit{crops} produits par YOLO
    sont conservés sous forme d'objets image en mémoire (PIL), évitant une
    relecture disque coûteuse avant l'OCR. Une légère marge de 3 pixels est
    ajoutée autour de chaque boîte pour ne pas tronquer les caractères de
    bord.
    \item \textbf{Seuil de confiance} ajustable (\texttt{conf} = 0,25 par
    défaut) pour filtrer les détections peu fiables.
\end{itemize}

% ─────────────────────────────────────────────────────────────────────────────
\subsection*{Reconnaissance du texte imprimé par TrOCR}

Une fois les zones \texttt{num\_compte} et \texttt{cmc7} localisées par
YOLOv8, leurs régions sont transmises au modèle \textbf{TrOCR} pour la
reconnaissance optique des caractères.

La variante retenue est \textbf{\texttt{microsoft/trocr-base-printed}},
spécialisée dans le texte \textit{imprimé} — par opposition à la variante
\textit{handwritten}. Ce choix est cohérent avec la nature des zones
traitées : le numéro de compte et la bande CMC7 sont des champs imprimés à
police régulière, pour lesquels un modèle entraîné sur du texte machine
offre la meilleure précision.

\paragraph{Optimisations OCR}

Deux optimisations majeures ont été appliquées pour accélérer l'OCR :

\begin{itemize}
    \item \textbf{Inférence par lot (batch)} : le numéro de compte et la
    bande CMC7 sont traités en \textbf{un seul passage} dans TrOCR plutôt
    qu'image par image, réduisant le surcoût d'appel du modèle.
    \item \textbf{Décodage glouton} (\textit{greedy decoding},
    \texttt{num\_beams = 1}) sous \texttt{torch.inference\_mode()}, qui
    désactive le calcul des gradients et privilégie la vitesse à
    l'exploration de la recherche par faisceau.
\end{itemize}

\paragraph{Post-traitement du numéro de compte}

Le texte brut renvoyé par TrOCR pour le numéro de compte est nettoyé : tous
les caractères non numériques sont supprimés afin de ne conserver que la
séquence de chiffres, garantissant un format exploitable en aval.

% ─────────────────────────────────────────────────────────────────────────────
\subsection*{Lecture du manuscrit par modèle Vision-Langage (VLM)}

\subsubsection{Motivation}

Les champs manuscrits du chèque — montant en chiffres, montant en lettres
et nom du bénéficiaire — constituent le cœur de la difficulté. L'écriture
manuscrite présente une variabilité considérable (style, inclinaison,
encre, ratures) que les approches OCR classiques, ainsi qu'un pipeline
YOLO + TrOCR, peinent à gérer de façon robuste.

Comme pour les justificatifs de domicile, un \textbf{modèle Vision-Langage
(VLM)} a été retenu pour cette tâche. Plutôt que de localiser puis lire
chaque champ, le VLM raisonne sur la sémantique globale du chèque et répond
directement à la question \textit{« quels sont les montants et le
bénéficiaire de ce chèque ? »}, en produisant un JSON structuré.

\subsubsection{Choix du modèle : Qwen2-VL-2B}

Le modèle retenu est \textbf{Qwen2-VL-2B-Instruct}, développé par Alibaba
Cloud. La variante 2B (2 milliards de paramètres) a été privilégiée pour :
\begin{itemize}
    \item ses capacités multilingues natives, incluant le français, langue
    des montants en lettres ;
    \item son architecture Vision-Langage unifiée, qui traite l'image et le
    texte dans un même espace de représentation ;
    \item son empreinte mémoire réduite, compatible avec un déploiement sur
    GPU à VRAM limitée.
\end{itemize}

Le \textit{prompt} fixe transmis au modèle est formulé en langage naturel :
\begin{quote}
\textit{« Lis ce chèque et renvoie UNIQUEMENT un JSON avec les clés
montant\_chiffre, montant\_lettre, beneficiaire. »}
\end{quote}

La réponse attendue est un JSON structuré directement exploitable :

\begin{verbatim}
{
  "montant_chiffre": "15000",
  "montant_lettre":  "QUINZE MILLE DIRHAMS",
  "beneficiaire":    "MOHAMMED ALAOUI"
}
\end{verbatim}

\subsubsection{Technique de fine-tuning : LoRA}

Le fine-tuning complet d'un modèle de 2 milliards de paramètres étant
coûteux et inutile pour une adaptation de domaine, la technique
\textbf{LoRA} (\textit{Low-Rank Adaptation}) \cite{hu2022lora} a été
employée. LoRA n'altère pas directement les poids du modèle pré-entraîné :
elle adjoint de petites matrices d'adaptation de rang réduit ($r = 16$) sur
les couches d'attention clés (\texttt{q\_proj}, \texttt{k\_proj},
\texttt{v\_proj}, \texttt{o\_proj}). Seuls ces adaptateurs sont mis à jour
pendant l'entraînement, ce qui réduit drastiquement la consommation mémoire
et le temps de calcul tout en préservant les connaissances générales du
modèle.

La quantisation \textbf{4-bit NF4} (NormalFloat4), appliquée via la
bibliothèque BitsAndBytes avec un \texttt{compute\_dtype} en bfloat16,
compresse les poids du modèle de base en mémoire sans dégradation notable
des performances.

\begin{table}[H]
\centering
\label{tab:vlm_cheque_config}
\begin{tabular}{|l|c|}
\hline
\textbf{Paramètre} & \textbf{Valeur} \\
\hline
Modèle de base            & Qwen2-VL-2B-Instruct          \\
Quantisation              & 4-bit NF4 (bfloat16)          \\
Technique d'adaptation    & LoRA ($r=16$, $\alpha=32$)    \\
Couches ciblées           & q\_proj, k\_proj, v\_proj, o\_proj \\
LoRA dropout              & 0,05                          \\
\hline
\end{tabular}
\caption{Configuration du fine-tuning du modèle Qwen2-VL-2B pour les chèques}
\end{table}

L'adaptateur final est stocké dans le dossier
\texttt{cheque\_qwen\_lora\_final/} et chargé à la volée par-dessus le
modèle de base au moment de l'inférence.

\subsubsection{Optimisation de l'inférence}

La génération du JSON s'effectue sous \texttt{torch.inference\_mode()} avec
un décodage glouton (\texttt{do\_sample = False}, \texttt{num\_beams = 1})
borné à 128 nouveaux tokens. La réponse textuelle est ensuite parsée par
une expression régulière qui extrait le bloc JSON, garantissant la
robustesse face à d'éventuels préambules ou commentaires générés par le
modèle.

% ─────────────────────────────────────────────────────────────────────────────
\subsection*{Validation métier}

Une fois les deux pipelines exécutés (manuscrit + zones), l'ensemble des
données extraites est confronté à un module de validation métier
(\texttt{tools/validator.py}) qui produit une décision structurée. Les
contrôles portent sur :
\begin{itemize}
    \item la présence du montant en chiffres et du montant en lettres ;
    \item la présence du nom du bénéficiaire ;
    \item la détection \textbf{et} la lisibilité du numéro de compte ;
    \item la détection de la bande CMC7 et sa lisibilité ;
    \item la présence de la signature.
\end{itemize}

La décision finale est déterminée par le nombre d'anomalies relevées :

\begin{table}[H]
\centering
\label{tab:validation_cheque}
\begin{tabular}{|l|l|}
\hline
\textbf{Statut} & \textbf{Condition} \\
\hline
\texttt{VALIDE}                  & Aucune anomalie : tous les champs présents et lisibles \\
\texttt{À VÉRIFIER MANUELLEMENT} & 1 ou 2 anomalies (contrôle opérateur) \\
\texttt{REJETÉ}                  & 3 anomalies ou plus, ou échec d'extraction en amont \\
\hline
\end{tabular}
\caption{Règles de décision de la validation métier du chèque}
\end{table}

Ce mécanisme s'inscrit dans la logique \textit{Human-in-the-Loop} du
système : les cas incertains ne sont jamais validés ni rejetés
automatiquement, mais routés vers un opérateur humain pour contrôle.

% ─────────────────────────────────────────────────────────────────────────────
\subsection*{Synthèse de l'architecture}

Le tableau ci-dessous récapitule la répartition des tâches entre les trois
modèles du pipeline chèque.

\begin{table}[H]
\centering
\label{tab:synthese_cheque}
\begin{tabular}{|p{3.0cm}|p{3.4cm}|p{4.2cm}|p{2.4cm}|}
\hline
\textbf{Modèle} & \textbf{Tâche} & \textbf{Champs} & \textbf{Type} \\
\hline
YOLOv8n &
Détection des zones &
num\_compte, signature, cmc7 &
Localisation \\
\hline
TrOCR (printed) &
OCR du texte imprimé &
num\_compte, cmc7 &
Reconnaissance \\
\hline
Qwen2-VL-2B + LoRA &
Lecture du manuscrit &
montant\_chiffre, montant\_lettre, beneficiaire &
Vision-Langage \\
\hline
\end{tabular}
\caption{Répartition des tâches entre les trois modèles du pipeline chèque}
\end{table}

Cette architecture hybride tire parti des forces de chaque famille de
modèles : la rapidité et la précision spatiale de YOLOv8 sur les zones
gabariées, l'efficacité de TrOCR sur le texte imprimé régulier, et la
capacité de raisonnement sémantique de Qwen2-VL sur l'écriture manuscrite.
L'ensemble est orchestré pour minimiser l'empreinte mémoire (modèles
chargés en singletons et libérables à la demande) et le temps d'inférence
(découpe en mémoire, OCR batché, décodage glouton sous
\texttt{inference\_mode}), garantissant un traitement de bout en bout fiable
et exploitable directement par le back-office.
